\section{Mechanism Evidence and Robustness}

\subsection{Mechanism evidence}

The central mechanism result is not that repair wins everywhere, but that
answer sensitivity, local process discrimination, and downstream utility are
linked in a stable way.

\begin{figure}[t]
\centering
\setlength{\fboxsep}{10pt}
\fbox{
\begin{minipage}{0.9\linewidth}
\small
\centering
\textbf{Higher local answer sensitivity}

\vspace{0.3em}
$\Downarrow$

\vspace{0.3em}
\textbf{Lower local \amcd}

\vspace{0.3em}
$\Downarrow$

\vspace{0.3em}
\begin{tabular}{c}
lower selection gain \\
higher exploitability
\end{tabular}

\vspace{0.75em}
\raggedright
Observed evidence:
\begin{itemize}[leftmargin=1.5em]
\item \texttt{corr(local\_ASS, local\_AMCD) = -0.2851}
\item \texttt{corr(local\_AMCD, selection\_gain\_at4) = 0.8475}
\item \texttt{corr(local\_AMCD, exploitability\_rate) = -0.6662}
\end{itemize}
\end{minipage}
}
\caption{Mechanism schematic. The key claim is not merely that
answer-sensitive behavior co-occurs with poor performance, but that the
observed correlations are consistent with lower local process discrimination
accompanying worse downstream utility.}
\label{fig:mechanism_chain}
\end{figure}

Across pooled head-level points:
\begin{itemize}[leftmargin=1.5em]
\item \texttt{corr(AMCD, selection\_gain\_at4) = 0.8778}
\item \texttt{corr(AMCD, exploitability\_rate) = -0.7724}
\item \texttt{corr(ordinary\_auroc, selection\_gain\_at4) = 0.4659}
\item \texttt{corr(ordinary\_auroc, exploitability\_rate) = -0.0104}
\end{itemize}
These head-level correlations already show that \amcd{} is much closer to
downstream utility than ordinary AUROC. Figure~\ref{fig:mechanism_chain}
summarizes the behavioral pathway most consistent with the pooled correlations
and the later robustness results.

The local analysis is stronger still:
\begin{itemize}[leftmargin=1.5em]
\item \texttt{corr(local\_ASS, local\_AMCD) = -0.2851}
\item \texttt{corr(local\_ASS, selection\_gain\_at4) = -0.2758}
\item \texttt{corr(local\_ASS, exploitability\_rate) = 0.2726}
\item \texttt{corr(local\_AMCD, selection\_gain\_at4) = 0.8475}
\item \texttt{corr(local\_AMCD, exploitability\_rate) = -0.6662}
\end{itemize}
The most consistent reading is that answer sensitivity is harmful not merely
because it exists, but because lower local process discrimination accompanies
it and tracks the downstream utility gap. This explains why ordinary AUROC can
remain misleadingly healthy while deployment behavior still degrades.

\subsection{Robustness}

The visible-vs-masked reranker gap also survives transfer checks.
On the naturalized deployment benchmark, \texttt{Qwen3-Reranker-8B} visible
falls from hard synthetic \amcd{} \texttt{0.6667} to structured generated
\texttt{0.3235}, then only recovers to \texttt{0.5294} on the naturalized
generated slice. The masked reranker instead goes from \texttt{0.7546} to
\texttt{0.8235} and then \texttt{0.8529}. The visible reranker is therefore
more brittle under generated OOD, while the masked reranker stays strong and
low-sensitivity.

Against the mixed visible-attacker ensemble, the strongest current target,
\texttt{reranker8\_masked}, remains robust:
\begin{itemize}[leftmargin=1.5em]
\item \texttt{pairwise\_attack\_win\_rate = 0.1944}
\item \texttt{quartet\_top\_attack\_rate = 0.1111}
\item attacked-quartet \selectiongain{} \texttt{= 0.3333}
\item attacked-quartet \exploitability{} \texttt{= 0.0833}
\end{itemize}
For comparison, \texttt{reranker8\_visible} under the same mixed attacker has
attacked-quartet \exploitability{} \texttt{= 0.25}, and
\texttt{visible\_cond\_swap} also remains at \texttt{0.25}. This is strong
evidence that masked deployment is not only better on average but also more
robust under transfer from visible-family failures.

The main remaining weak slice for the best current model is planning-style
naturalized blocksworld. On that slice, \texttt{reranker8\_masked} reaches
\auroc{} \texttt{= 0.6875}, \amcd{} \texttt{= 0.5833}, and \asstotal{}
\texttt{= 0.0075}. The remaining weakness is therefore concrete rather than
generic. The cleaner \textsc{Craft-Clean} regime sharpens a complementary
boundary: once obvious artifact effects are reduced, the remaining weakness is
less about blocksworld surface artifacts and more about algebra and graph-path
faithfulness.
